\documentclass[11pt,letterpaper]{article}
\usepackage[margin=1in]{geometry}
\usepackage{graphicx} % Required for inserting images
\usepackage{amsthm,amsmath,amssymb}
\usepackage{algorithm}
\usepackage{algpseudocode}
\usepackage{xcolor}
\usepackage{soul}
\usepackage{enumitem}

% Theorems
\newtheorem{theorem}{Theorem}[section]
\newtheorem{corollary}[theorem]{Corollary}
\newtheorem{lemma}[theorem]{Lemma}
\newtheorem{claim}[theorem]{Claim}
\theoremstyle{definition}
\newtheorem{definition}[theorem]{Definition}
\newtheorem{remark}[theorem]{Remark}
\newtheorem{observation}[theorem]{Observation}

\newcommand{\reals}{\mathbb{R}}
\newcommand{\Span}{\mathrm{span}}
\newcommand{\coreq}{\mathrm{coreq}}

\newcommand{\calA}{\mathcal{A}}
\newcommand{\calB}{\mathcal{B}}
\newcommand{\calC}{\mathcal{C}}
\newcommand{\calD}{\mathcal{D}}
\newcommand{\calE}{\mathcal{E}}
\newcommand{\calF}{\mathcal{F}}
\newcommand{\calG}{\mathcal{G}}
\newcommand{\calH}{\mathcal{H}}
\newcommand{\calI}{\mathcal{I}}
\newcommand{\calJ}{\mathcal{J}}
\newcommand{\calK}{\mathcal{K}}
\newcommand{\calL}{\mathcal{L}}
\newcommand{\calM}{\mathcal{M}}
\newcommand{\calN}{\mathcal{N}}
\newcommand{\calO}{\mathcal{O}}
\newcommand{\calP}{\mathcal{P}}
\newcommand{\calQ}{\mathcal{Q}}
\newcommand{\calR}{\mathcal{R}}
\newcommand{\calS}{\mathcal{S}}
\newcommand{\calT}{\mathcal{T}}
\newcommand{\calU}{\mathcal{U}}
\newcommand{\calV}{\mathcal{V}}
\newcommand{\calW}{\mathcal{W}}
\newcommand{\calX}{\mathcal{X}}
\newcommand{\calY}{\mathcal{Y}}
\newcommand{\calZ}{\mathcal{Z}}

\newcommand{\cost}{\mathrm{COST}}
\newcommand{\ALG}{\mathrm{ALG}}
\newcommand{\OPT}{\mathrm{OPT}}
\newcommand{\Flow}{\textsc{Flow}}
\newcommand{\ouralg}{\textsc{LineTracker}}
\newcommand{\interval}[1]{\lfloor #1 \rceil}
\newcommand{\off}{\mathrm{off}}
\newcommand{\on}{\mathrm{on}}
\newcommand{\org}{\mathrm{org}}
\newcommand{\rel}{\mathrm{rel}}
\newcommand{\dis}[1]{#1}
\newcommand{\new}{\mathrm{new}}

\newcommand{\solution}{\medskip\noindent\textbf{Solution.}\par\medskip}

\usepackage[hidelinks]{hyperref}
\usepackage{cleveref}


\title{47-835: Network Optimization I \\ Homework 2}
\author{Poon Tze-Yang}
\date{September 2026}

 
\begin{document}

\maketitle

Students in the OR and ACO programs should answer all questions. For the others, the questions marked * are optional. Those marked ** are optional for all.

\begin{enumerate}
\item Consider the following ``coloring'' algorithm. Initially, all edges are uncolored. Then apply the following rules in arbitrary order until all edges are colored.

\noindent \textbf{Blue rule:} Select a cut containing no blue edge. Among the uncolored edges in the cut, select one of minimum cost and color it blue.

\noindent \textbf{Red rule:} Select a circuit (cycle) containing no red edge. Among the uncolored edges in the circuit, select one of maximum cost and color it red.
\begin{enumerate}
\item Show that as long as there are uncolored edges left, one of the rules is applicable.
\item Show that at all points, there exists a minimum spanning tree containing all blue edges but no red edge, so in the end the algorithm solves the MST problem.
\item Show that Kruskal's and Prim's algorithms are special cases of this algorithm.
\end{enumerate}
\end{enumerate}

\solution
 
\paragraph{(a)}
Suppose there exists an uncolored edge $e=uv$. 
Let $V_u$ and $V_v$ be the vertices reachable from $u$ and $v$, respectively, in the blue subgraph.
If $V_u = V_v$, there exists a blue path $P$ from $u$ to $v$. 
But $P+u$ forms a cycle with no red edges, which means the red rule can be applied.
Otherwise, the cut $\delta(V_u)$ forms a cut with no blue edges, and thus the blue rule can be applied.
\qed

\paragraph{(b)}
On a partially colored graph, let $p$ be the cost of an MST $T$ which uses all blue edges and possibly some uncolored edges.
We will consider the next edge to be colored $e=uv$ and let $p'$ be the same quantity on the new graph.
We show that both cases of the color of $e$, we have $p=p'$.

If $e$ is colored blue and $e\in T$, we are done since $T$ is still an MST in the new graph that uses all blue edges. 
If $e\notin T$, consider the fundamental cycle $C$ of $e$. 
Furthermore, from the cut defined by the blue rule, there must exist an uncolored $e'\in (C - e)$ which also crosses the cut such that $c(e')\geq c(e)$.
Thus, $c(T-e'+e) \leq c(T)$, and $T-e' + e$ forms our desired MST in the new graph.

Otherwise, $e$ is colored red. 
If $e\notin T$, we are done, since $T$ is still in the new graph.
Otherwise, $e\in T$, and let $C$ be the cycle used by the red rule. 
Deleting $e$ from $T$ splits the remaining edges into two connected components, and there exists some $e'\in C\backslash T$ which connects these two components.
By the red rule, $c(e)\geq c(e')$.
Thus, $c(T-e+e')\leq c(T)$, and $T-e+e'$ forms our desired MST in the new graph.
\qed

\paragraph{(c)}
Kruskal's repeatedly applies the blue rule with the connected blue component incident to the cheapest edge across components.
Prim's repeatedly applies the blue rule, always growing the same connected component.


\begin{enumerate}[resume]
\item Alice and Bob play the following game alternatively on a given simple connected graph $G$: Alice starts the game. At each of her turns, she deletes an edge of $G$. Following Alice, Bob fixes an edge at each of his turns. A fixed edge can not be deleted. Bob wins if he builds a spanning tree of $G$, otherwise, Alice wins. Show that if $G$ has two edge-disjoint spanning trees, then Bob wins the game (In fact, this is also the necessary condition).
\footnote{Serial exchange on the order of selection gives this.}
\end{enumerate}

\solution
We describe and prove a winning strategy for Bob.
Initialize the two spanning trees as the edge-disjoint trees $T_1$ and $T_2$ and let the set of edges Bob takes as $B\gets \emptyset$.
As the game proceeds we evolve $T_1$, $T_2$ and $B$ such that at the end of each of Bob's turns, we maintain the following invariants:
\begin{enumerate}
    \item $T_1$ and $T_2$ remain spanning trees;
    \item $T_1\cap T_2 = B$.
\end{enumerate}
Thus, at the game ends when $|B|=n-1$ and $B=T_1=T_2$ is a spanning tree.

At the start of the $i$th round, we consider two cases on the edge $e$ which Alice deletes:

\textbf{Case 1: }$e\in T_1\triangle T_2$. Without loss of generality, suppose $e\in T_1$. 
Let $C$ be the fundamental cycle of $e$ in $T_2$.
Then deleting $e$ from $T_1$ produces two connected components $D_1$ and $D_2$.
There must exist some edge $e'\in C$ connecting $D_1$ and $D_2$.
Furthermore, since $e'\notin T_1$, therefore $e'\in T_2\backslash T_1$.
Thus, Bob can select $e'$, so $B\gets B+e'$. 
The invariants are maintained, since $T_1\gets T_1-e+e'$ remains a spanning tree, $T_2$ is unchanged and $T_1\cap T_2=B$.

\textbf{Case 2: }$e\notin T_1\cup T_2$. 
Bob can select any edge $e\in T_1$.
Let $C$ be the fundamental cycle of $e$ in $T_2$, and select any $e'\in (C-e)$.
The invariants are maintained since $T_1$ is unchanged, $T_2\gets T_2-e'+e$ remains a spanning tree and $T_1\cap T_2=B$.

\qed

\begin{enumerate}[resume]
\item Given a graph $G$, we start modifying it as follows: In each round, as long as $G$ is cyclic, pick an arbitrary cycle of $G$, e.g. $C$, and remove all the edges of $C$ from $G$. Then, add a new vertex to the graph, e.g. $v_C$, and connect $v_C$ to all the vertices of $C$ (note that the newly added vertex can be on the cycles that we pick in the next rounds). Show that regardless of the choice of cycles, $G$ becomes a forest (collection of trees) after a finite number of rounds.
\end{enumerate}

\solution
(A counting proof)

Note that after the operation, the total number of vertices in the graph increases by $1$, but the number of edges remains the same. 
Furthermore, all vertices after each operation have degree at least 1: we only delete edges from vertices of $C$, but we also add one edge to each such vertex.

Hence, if the starting number of edges is $k$, after at most $k$ iterations all vertices have degree at most $1$. 
Therefore, there are no remaining cycles, since vertices of a cycle have degree at least two.


\begin{enumerate}[resume]
\item The tree graph of a connected graph $G$ is the graph whose vertices are the spanning trees $T_1, T_2, \ldots, T_r$ of $G$, with an edge between $T_i$ and $T_j$ if and only if they have exactly $n-2$ edges in common. Show that the tree graph of any connected graph is connected.
\footnote{Serial exchange works for this too.}
\end{enumerate}

\solution
By the hypothesis, our graph is connected, so we have at least one spanning tree.
If spanning trees $T_i$ and $T_j$ have exactly $n-2$ edges in common, one edge swap transforms one into the other.

As shown in lectures, any spanning tree can be transformed into a minimum spanning tree by a sequence of swaps. 
By assigning weights that increase by a factor of two to the edges, we have a unique minimum spanning tree.
Therefore, all trees must have a path to the minimum spanning tree in the tree graph, so all trees are connected in the tree graph.

\begin{enumerate}[resume]
\item Recall that an edge cut is the set of all edges going from a subset of nodes $S$ to its complement $V \setminus S$, and denoted $[S, V \setminus S]$. Prove that the symmetric difference of two different edge cuts is an edge cut.
\end{enumerate}

\solution
Let the two cuts be defined by vertex sets $P,Q$. 
This partitions the vertex set into four (possibly empty) groups:
\begin{align*}
    A:= P\cap Q, \quad B:=P\cap \overline{Q}, \quad C:=\overline{P}\cap Q, \quad D:=\overline{P}\cap \overline{Q}.
\end{align*}
Let $E_{XY}$ be the edges between components $X$ and $Y$. We have:
\begin{align*}
    \delta(P) = E_{AC} \cup E_{BD}\cup E_{AD}\cup E_{BC}, \\
    \delta(Q) = E_{AB} \cup E_{CD}\cup E_{AD}\cup E_{BC}. 
\end{align*}
So 
\begin{align*}
    \delta(P) \triangle \delta(Q) = E_{AB} \cup E_{CD} \cup E_{AC} \cup E_{BD} = \delta(A\cup D),
\end{align*}
which is a cut.

\begin{enumerate}[resume]
\item Given an $n$-node directed acyclic graph (DAG), show that the nodes can be numbered with distinct integers from $1$ to $n$ such that all arcs go from smaller to higher numbered nodes. This numbering is called a topological ordering of the graph, and is commonly used in computing quantities on the nodes via dynamic programming.
\end{enumerate}

\solution

\begin{lemma}
    For any vertex $u$ in a DAG $G$, let $R_u\subseteq V$ be the set of vertices reachable from $u$. 
    Then (i) all vertices of $V\backslash R_u$ are not reachable in $G$ from any vertex of $R_u$; (ii) $u$ is not reachable from any vertex in $R_u - u$; and (iii) $G\backslash R_u$ is also a DAG.
\end{lemma}
\begin{proof}
    (i): Clearly, otherwise that vertex would also be in $R_u$.
    
    (ii): Any $v\in R_u$ is reachable from $u$, so if there exists a $(v,u)$-path, then there exists a cycle in $G$.

    (iii): Any cycle in $G\backslash R_u$ is also a cycle in $G$.
\end{proof}

Our algorithm proceeds inductively on the size of $G$.
If $|V|=1$, the topological ordering is just one vertex.
Otherwise, select any $u\in V$ and define $A = V\backslash R_u$ and $B = R_u - u$. 

\begin{lemma}
    There exists a valid topological ordering beginning with the vertices of $A$ in some order, then $u$, then the vertices of $B$ in some order.
\end{lemma}
\begin{proof}
    By the inductive hypothesis, there is a valid topological ordering of the vertices of $A$ and $B$.
    Our ordering will be that of $A$, then $u$, then that of $B$.
    Call an edge \emph{reverse-order} if it is directed from vertex later in the ordering to a vertex earlier in the ordering.

    We show that there are no edges of $G$ crossing the partition $A\sqcup \{u\} \sqcup B$ in the reverse order. 
    There are no reverse-order edges within $A$ or $B$ because they are alredy topologically ordered.
    There is no reverse-order edge $(x,y)$ ending in $A$ because then $y$ should be in $R_u$.
    There is no reverse-order edge from $B$ to $u$ because that would be a cycle in $G$.
\end{proof}

Thus, our algorithm can recursively topologically order $A$ and $B$, giving a topological ordering of $G$.


\begin{enumerate}[resume]
\item[7.*] An edge $e$ is a cut edge if $G - e$ has one more connected component than $G$ (in undirected graphs). Show that
\begin{enumerate}
\item if each vertex degree in $G$ is even, then $G$ has no cut edge.
\item if $G$ is a $k$-regular bipartite graph with $k \geq 2$, then $G$ has no cut edge.
\end{enumerate}
\end{enumerate}

\solution
\paragraph{(a)}
\begin{claim}
    Any connected component has an even total degree.
\end{claim}
\begin{proof}
    The sum of all degrees is equal to twice the number of edges.
\end{proof}
Suppose for contradiction that there exists a cut edge $e$, which separates component $C$ into $C_1$ and $C_2$.
Deleting $e$ decreases the total degree of $V(C_1)$ by $1$. Since all vertices were originally of even degree, this means $C_1$ has odd total degree, a contradiction.

\paragraph{(b)}
Suppose for contradiction that there is a cut edge $e= uv$. 
Restrict our attention to the connected component $C$ of $e$ in $G$.
Furthermore, let $C_u$ be the connected component of $G-e$ containing $u$, and let $C_v$ be the connected component of $G-e$ containing $v$.
Then all vertices of $C_u$ are of degree $k$ except for $u$, which is of degree $k-1$.

Let the bipartition of $V(C_u)$ be $A\sqcup B$, and without loss of generality let $u\in A$.
Let $a:=|A|$ and $b:=|B|$. 
Then by counting the degrees of each side,
\begin{align*}
    &(a-1)k + k-1 = bk \\
    &\Rightarrow (a-b)k = 1.
\end{align*}
But $a,b$ are integer and $k\geq 2$, which is a contradiction.

\begin{enumerate}[resume]
\item[8.*] Let $G$ be a strongly connected graph with $2n-1$ (directed) edges, i.e., for any ordered pair of vertices $u, v$, there is a directed path from $u$ to $v$. Prove that there exist an edge $e$ in $G$ such that $G - e$ is still strongly connected.
\end{enumerate}

\solution
Select any vertex $r\in G$. 
Since $G$ is strongly connected, there exists an $r$-arborescence $B$ rooted at $r$.
The arborescence $B$ has $n-1$ arcs, so there are $n$ arcs of $G$ remaining.
Since $G$ is strongly connected, each vertex $u\neq r$ has some path to $r$.

If $r$ has an arc outside of $B$, removing it preserves strong connectivity, since $r$ can still reach all other vertices using arcs in $B$ and any other vertex would not use that arc on its shortest path to $r$.
Otherwise, no arcs outside $B$ start at $r$.

For each vertex $u\neq r$, mark one arc starting at $u$ which is on a shortest path from $u$ to $r$. 
If there are multiple such arcs for one $u$, mark any of them. 
Let the set of marked arcs be $C$.

\begin{lemma}
    The graph with arcs $B\cup C$ is strongly connected.
\end{lemma}
\begin{proof}
    Since we have all arcs of $B$, $r$ can reach all other vertices. 
    Every other vertex can reach $r$ by following the arcs of $C$: for vertex $u\neq r$ with arc $uv\in C$, $v$ is one arc closer to $C$ than $u$.
    Hence, all vertices can reach any other vertex by first going through $r$. 
\end{proof}

There are $n-1$ marked arcs in $C$, since there are $n-1$ vertices in $V-r$. 
Hence, $|B\cup C| \leq |B| + |C| = 2n-2$. 
Therefore, there is at least one arc of $G$ that can be removed while preserving strong connectivity.

\begin{enumerate}[resume]
\item[9.*] Show that the spanning tree polyhedron of Edmonds with `subtour elimination' constraints describing the convex hull of spanning trees of a connected undirected graph $(V,E)$ has integral solutions whenever the edge costs are nonnegative integers.
\begin{align*}
    x(E(S)) &\leq |S| - 1 && \forall S \subset V \\
    x(E(V)) &= |V| - 1 \\
    x_e &\geq 0 && \forall e \in E
\end{align*}
Hint: Write the dual out carefully and give appropriate values to these duals based on the minimum spanning tree found by Kruskal's greedy algorithm.
\end{enumerate}

\solution
% Subtracting each $S$ constraint from the $V$ constraint gives the equivalent formulation:

% \begin{align*}
%     \min \sum_{e} c_e x_e \\
%     x(E(V\backslash S)) &\geq |V\backslash S| && \forall S \subset V \\
%     x(E(V)) &= |V| - 1 \\
%     x_e &\geq 0 && \forall e \in E
% \end{align*}

% This gives the following dual

% \begin{align*}
%     \max \sum_{S\subset V} (|S|-1)y_S\\
%     \sum_{S:e\notin E(S)} y_S\leq c_e && \forall e\in E \\
%     y_S\geq 0 && \forall S\subset V 
% \end{align*}

For each $S\subseteq V$, we have the dual variable $y_s$.
Assume the primal problem has a non-negative, integral cost vector $c$.
The dual is as follows:

\begin{align*}
    \max \ (|V|-1)y_V +\sum_{S\subseteq V} (|S|-1)y_S\\
    y_V + \sum_{S\subset V:e\in E(S)} y_S\leq c_e && \forall e\in E \\
    y_S\leq 0 && \forall S\subset V 
\end{align*}

Consider Kruskal's algorithm on the given graph. 
Let the sequence of edges it selects be $E^\ALG = \langle e_1, .. , e_{n-1}\rangle$.
This produces a sequence of forests $F_i:=\bigcup_{j=1}^i \{e_j\}$ for $i\in \{0 .. n-1\}$, where $F_0=\emptyset$ and $F_{n-1}$ is a minimum spanning tree.
For each $i \in \{0.. n-2\}$, let $C_i$ be the connected-component containing $e_i$ in $F_i$, and let $e'_i$ be next edge added to $C_i$.
Let the collection of these components and the final tree $F_{n-1}$ be $\calC$.

Consider the solution $y$ such that 

\begin{align*}
    y_S = \begin{cases}
        c_{e_{n-1}} & \text{for } S=V, \\
        -(c_{e'_i} - c_{e_i})  & \text{for } S = V(C_i), \\
        0 & \text{otherwise.}
    \end{cases}
\end{align*}

\begin{lemma}
    $y$ is dual-feasible.
\end{lemma}
\begin{proof}
    Since Kruskal's adds each $e'_i$ at a later time than $e_i$ we have that $c_{e'_i}\geq c_{e_i}$, so all non-positivity constraints are satisfied.
    For each $e\in E$, consider the following cases
    \begin{enumerate}
        \item \ul{$e \in E^\ALG$.} 
        The only non-zero dual variables in the corresponding constraints are the elements of $\calC$ containing $e$.
        Let $E'':=\langle e''_0=e, .. , e''_k = e_{n-1}\rangle$ be the sequence of edges starting at $e$ such that $e''_{j+1} = e'_j$, i.e., $e''_{j+1}$ is the next edge added to the first component containing $e_j$.
        Note that this sequence must end in $e_{n-1}$, since $e_{n-1}$ is the last edge added and there is only one component in the final spanning tree.
        Hence, we get
        \begin{align*}
            y_V + \sum_{S\subset V:e\in E(S)} y_S & =\sum_{C\in \calC: e\in C} y_{V(C)} \\
            & = c_{e_{n-1}} + \sum_{j=0}^{k-1} \left(- c_{e''_{j+1}} + e''_j\right) \\
            & = c_{e_{n-1}} - c_{e_{n-1}} + c_{e}&&(\text{telescoping}) \\
            & = c_e. 
        \end{align*}
        Thus, the corresponding constraint holds.
        \item \ul{$e \notin E^\ALG$.} Let $C_i$ be the first component that contains $e$, such that $e_i$ is the corresponding edge in $E^\ALG$ that created that component.
        It follows that the terms on the left-hand-side of the constraint corresponding to $e$ are the same as those of $e_i$'s, since $e$ and $e_i$ share the same first component.
        
        Additionally, $e_i$ and $e$ sharing the first component means that at the time when $e_i$ was first picked, $e$ was not spanned by any existing component.
        According to Kruskal's this means that $c_e\geq c_{e_i}$. Hence, the constraint corresponding to $e$ also holds.
    \end{enumerate}
\end{proof}

Next, we show that the dual objective value attained by $y$ is equal to the primal objective value attained by Kruskal's. 
Together with the fact that all $y_S$ values are integer if all $c_e$ are integer non-negative, This shows that the primal optimal value is also integer, completing the proof.

\begin{lemma}
    The dual objective value attained by $y$ is equal to the primal objective value attained by Kruskal's.
\end{lemma}
\begin{proof}
    (Sketch): Since each component is a tree, the coefficient exactly splits $y_S$ of value into each edge of the final tree it contains.
    Then, we can redistribute the objective such that each edge taken by Kruskal's gets exactly $c_e$ of value.
\end{proof}


\begin{enumerate}[resume]
\item[10.**] Given a complete graph on nodes $V$ with a cost $c_{uv}$ between every pair of nodes $\{u,v\}$, partition $V$ into two parts $V_1, V_2$ such that the minimum cost of any edge between nodes in the same part is maximized. In other words, we want to maximize
\[
\min\left\{ \min_{u,v \in V_1} c_{uv}, \min_{u,v \in V_2} c_{uv} \right\}.
\]
Give an efficient algorithm for the task. Show an easy $O(|V|)$ procedure to check whether the solution is unique.
\end{enumerate}

\solution

Algorithm: 
Starting from an empty graph, add edges in ascending order of weight.
After adding each edge, check whether the graph contains an odd cycle.
We can compute this by trying to 2-color the graph.
Return the bipartition just before we get an odd cycle.

Runtime: 
2-coloring can be done by bfs.
At most $E$ iterations.

Correctness:
Any bipartition of the vertices of the odd cycle must contain a monochromatic edge of the cycle.
Hence, we canot do better than cutting every edge of the odd cycle less the most expensive edge.

Unique solution:
For each connected component of a graph with a valid 2-coloring, 

\end{document}
